% ═══════════════════════════════════════════════════════════════
% LEHRERHINWEISE: Las direcciones (Wegbeschreibung)
% ═══════════════════════════════════════════════════════════════
% Sequenz 6c: Mi ciudad - Las direcciones
% Fachfarbe: #c41e3a (Karminrot)
% ═══════════════════════════════════════════════════════════════

\documentclass[11pt, a4paper]{article}

% ═══════════════════════════════════════════════════════════════
% SW-MODUS TOGGLE
% ═══════════════════════════════════════════════════════════════
\newif\ifswmodus
\swmodusfalse

% ═══════════════════════════════════════════════════════════════
% PAKETE
% ═══════════════════════════════════════════════════════════════

\usepackage[utf8]{inputenc}
\usepackage[T1]{fontenc}
\usepackage[ngerman]{babel}
\usepackage{helvet}
\renewcommand{\familydefault}{\sfdefault}

\usepackage[margin=2cm]{geometry}
\usepackage{enumitem}
\usepackage{tabularx}
\usepackage{booktabs}
\usepackage{longtable}

\usepackage{xcolor}
\usepackage{amssymb}
\usepackage{tcolorbox}
\tcbuselibrary{skins}

\usepackage{fancyhdr}
\usepackage{lastpage}
\usepackage{needspace}

% ═══════════════════════════════════════════════════════════════
% FARBEN
% ═══════════════════════════════════════════════════════════════

\definecolor{spanisch-color}{HTML}{c41e3a}
\definecolor{grau-color}{HTML}{888888}
\definecolor{hellgrau-color}{HTML}{f5f5f5}
\definecolor{basis-color}{HTML}{2a7a4b}
\definecolor{standard-color}{HTML}{2c5aa0}
\definecolor{erweiterung-color}{HTML}{7b1fa2}
\definecolor{loesung-color}{HTML}{c62828}

\definecolor{sw-dark}{HTML}{000000}
\definecolor{sw-gray}{HTML}{666666}
\definecolor{sw-white}{HTML}{ffffff}

\ifswmodus
    \colorlet{spanisch}{sw-dark}
    \colorlet{grau}{sw-gray}
    \colorlet{hellgrau}{sw-white}
    \colorlet{basis}{sw-dark}
    \colorlet{standard}{sw-dark}
    \colorlet{erweiterung}{sw-dark}
    \colorlet{loesung}{sw-dark}
\else
    \colorlet{spanisch}{spanisch-color}
    \colorlet{grau}{grau-color}
    \colorlet{hellgrau}{hellgrau-color}
    \colorlet{basis}{basis-color}
    \colorlet{standard}{standard-color}
    \colorlet{erweiterung}{erweiterung-color}
    \colorlet{loesung}{loesung-color}
\fi

\newcommand{\fachfarbe}{spanisch}

% ═══════════════════════════════════════════════════════════════
% VARIABLEN
% ═══════════════════════════════════════════════════════════════

\newcommand{\thema}{Las direcciones -- Wegbeschreibung}
\newcommand{\sequenz}{06}
\newcommand{\block}{c}
\newcommand{\fach}{Spanisch}
\newcommand{\klasse}{7}
\newcommand{\dauer}{90 Min. (Doppelstunde)}
\newcommand{\lerngruppengroesse}{24}
\newcommand{\datum}{\today}

% ═══════════════════════════════════════════════════════════════
% KOPF- UND FUSSZEILE
% ═══════════════════════════════════════════════════════════════

\pagestyle{fancy}
\fancyhf{}
\renewcommand{\headrulewidth}{0.4pt}
\renewcommand{\footrulewidth}{0.4pt}

\lhead{\fach{} -- Klasse \klasse}
\chead{\textbf{LH-\sequenz\block{} (Lehrerhinweise)}}
\rhead{\datum}

\lfoot{}
\cfoot{}
\rfoot{Seite \thepage\ von \pageref{LastPage}}

% ═══════════════════════════════════════════════════════════════
% BEFEHLE
% ═══════════════════════════════════════════════════════════════

\newcommand{\B}{\textcolor{basis}{\textbf{[B]}}}
\newcommand{\Std}{\textcolor{standard}{\textbf{[S]}}}
\newcommand{\E}{\textcolor{erweiterung}{\textbf{[E]}}}

\newcommand{\loesung}[1]{\textcolor{loesung}{\textbf{#1}}}

\newcommand{\es}[1]{\textcolor{\fachfarbe}{\textbf{#1}}}

\newtcolorbox{kurzplan}{
    colback=hellgrau,
    colframe=\fachfarbe,
    colbacktitle=white,
    coltitle=\fachfarbe,
    boxrule=1pt,
    arc=0pt,
    fonttitle=\bfseries,
    attach boxed title to top left={yshift=-2mm, xshift=5mm},
    boxed title style={boxrule=0pt, colback=white},
    title=Kurzplan
}

\newtcolorbox{differenzierung}{
    colback=white,
    colframe=grau,
    colbacktitle=white,
    coltitle=grau,
    boxrule=0.5pt,
    arc=0pt,
    fonttitle=\bfseries,
    attach boxed title to top left={yshift=-2mm, xshift=5mm},
    boxed title style={boxrule=0pt, colback=white},
    title=Differenzierung
}

\newtcolorbox{fehlerbox}{
    colback=white,
    colframe=loesung,
    colbacktitle=white,
    coltitle=loesung,
    boxrule=0.5pt,
    arc=0pt,
    fonttitle=\bfseries,
    attach boxed title to top left={yshift=-2mm, xshift=5mm},
    boxed title style={boxrule=0pt, colback=white},
    title=Häufige Fehler
}

\newtcolorbox{materialbox}{
    colback=white,
    colframe=\fachfarbe,
    colbacktitle=white,
    coltitle=\fachfarbe,
    boxrule=0.5pt,
    arc=0pt,
    fonttitle=\bfseries,
    attach boxed title to top left={yshift=-2mm, xshift=5mm},
    boxed title style={boxrule=0pt, colback=white},
    title=Material-Checkliste
}

% ═══════════════════════════════════════════════════════════════
% DOKUMENT
% ═══════════════════════════════════════════════════════════════

\begin{document}

% ═══════════════════════════════════════════════════════════════
% KOPFBEREICH
% ═══════════════════════════════════════════════════════════════

\begin{center}
    {\Large\bfseries\textcolor{\fachfarbe}{Lehrerhinweise: \thema}}\\[0.3em]
    {\normalsize LH-\sequenz\block{} | \dauer}
\end{center}

\vspace{10pt}

% ═══════════════════════════════════════════════════════════════
% 1. KURZPLAN
% ═══════════════════════════════════════════════════════════════

\section*{1. Stundenverlauf (Kurzplan)}

\textbf{1. Stunde (45 Min.)}

\begin{tabularx}{\textwidth}{|c|l|X|l|}
\hline
\textbf{Zeit} & \textbf{Phase} & \textbf{Inhalt} & \textbf{Material} \\
\hline
0--5 & Einstieg & Stadtplan zeigen, Frage: \es{¿Cómo llego a...?} & Beamer/Tafel \\
\hline
5--15 & Input & Richtungsangaben mit Gesten einführen & Bildkarten \\
\hline
15--25 & Übung 1 & TPR: Anweisungen befolgen (todo recto, gira...) & -- \\
\hline
25--35 & Erweiterung & Ortsangaben: al lado de, enfrente de, cerca/lejos & Bildkarten \\
\hline
35--45 & Sicherung & Vokabelabfrage, Zusammenfassung & Tafel \\
\hline
\end{tabularx}

\vspace{1em}

\textbf{2. Stunde (45 Min.)}

\begin{tabularx}{\textwidth}{|c|l|X|l|}
\hline
\textbf{Zeit} & \textbf{Phase} & \textbf{Inhalt} & \textbf{Material} \\
\hline
0--5 & Wiederholung & Schnellquiz: Pfeile zeigen, Vokabeln nennen & Bildkarten \\
\hline
5--10 & Modell & Beispieldialog vorführen (L + Schüler) & -- \\
\hline
10--25 & Anwendung & Rollenspiel: Tourist fragt nach dem Weg & Rollenkarten \\
\hline
25--35 & Festigung & AB-06c bearbeiten (Aufgaben 1-4) & AB \\
\hline
35--45 & Abschluss & Lösungen besprechen, HA erklären & Tafel \\
\hline
\end{tabularx}

% ═══════════════════════════════════════════════════════════════
% 2. DIFFERENZIERUNG
% ═══════════════════════════════════════════════════════════════

\section*{2. Differenzierung}

\begin{differenzierung}
\textbf{Legende:} \B{} Basis (BR) | \Std{} Standard (MR) | \E{} Erweiterung (GY)

\vspace{0.5em}

\textbf{WICHTIG:} Diese Marker sind NUR für die Lehrkraft. Auf Schülermaterial erscheinen KEINE Niveaumarkierungen!
\end{differenzierung}

\vspace{1em}

\begin{tabularx}{\textwidth}{|l|X|X|X|}
\hline
& \B{} \textbf{Basis} & \Std{} \textbf{Standard} & \E{} \textbf{Erweiterung} \\
\hline
\textbf{Aufgabe 1} & Mit Bildunterstützung & Selbstständig & + Sätze bilden \\
\hline
\textbf{Aufgabe 2} & Wortliste nutzen & Selbstständig & + Eigene Sätze \\
\hline
\textbf{Aufgabe 3} & Satzbausteine erhalten & Freie Formulierung & Komplexer Dialog \\
\hline
\textbf{Aufgabe 4} & 2 einfache Sätze & 3 Sätze wie vorgegeben & 4+ Sätze, Alternativen \\
\hline
\textbf{Rollenspiel} & Vorgegebener Dialog ablesen & Dialog mit Stichwörtern & Freies Rollenspiel \\
\hline
\end{tabularx}

\vspace{1em}

\textbf{Konkrete Hinweise:}
\begin{itemize}
    \item \B{} SuS mit LRS: Vergrößerte AB-Version, Zeitzugabe (+10 Min.)
    \item \B{} DaZ-SuS: Wortschatz deutsch-spanisch vorab erklären, Bildkarten nutzen
    \item \E{} Leistungsstarke: Rollenspiel mit Komplikationen (z.B. ``Die Straße ist gesperrt'')
\end{itemize}

% ═══════════════════════════════════════════════════════════════
% 3. HÄUFIGE FEHLER
% ═══════════════════════════════════════════════════════════════

\section*{3. Häufige Fehler}

\begin{fehlerbox}
\begin{tabularx}{\textwidth}{|l|X|X|}
\hline
\textbf{Fehler} & \textbf{Beispiel} & \textbf{Reaktion} \\
\hline
Verwechslung rechts/links & *\textit{a la derecha} statt \es{a la izquierda} & Gesten etablieren; Eselsbrücke: ``i'' in izquierda = links \\
\hline
Aussprache \textit{derecha} & *[derecha] statt [de'recha] & Chorisches Sprechen, /ch/ wie ``tsch'' betonen \\
\hline
Präposition vergessen & *\textit{al lado el banco} statt \es{al lado del banco} & Als feste Wendung einüben: ``al lado de'' \\
\hline
Artikel + de kontrahiert & *\textit{de el} statt \es{del} & Regel: de + el = del (immer!) \\
\hline
Deutsche Struktur & *\textit{La estación es lejos} statt \es{está lejos} & ser vs. estar: Ort = estar \\
\hline
\end{tabularx}
\end{fehlerbox}

% ═══════════════════════════════════════════════════════════════
% 4. MATERIAL-CHECKLISTE
% ═══════════════════════════════════════════════════════════════

\section*{4. Material-Checkliste}

\begin{materialbox}
\begin{itemize}[label=$\square$]
    \item AB-\sequenz\block.pdf (\lerngruppengroesse{} Kopien)
    \item ML-\sequenz\block.pdf (1× für Lehrkraft)
    \item Lehrwerk: Qué pasa 1, S. 104--108
    \item Bildkarten: Richtungspfeile (todo recto, derecha, izquierda)
    \item Stadtplan (vergrößert für Tafel/Beamer)
    \item Rollenkarten: Tourist/Einheimischer (12 Sets)
    \item Tafelmarker / Kreide
    \item Optional: LP-\sequenz\block.html (Tablets/PCs)
\end{itemize}
\end{materialbox}

% ═══════════════════════════════════════════════════════════════
% 5. LÖSUNGEN
% ═══════════════════════════════════════════════════════════════

\section*{5. Lösungen AB-\sequenz\block}

\textbf{Merkkasten:}

\begin{tabular}{ll}
\es{todo recto} & $\rightarrow$ \loesung{geradeaus} \\
\es{a la derecha} & $\rightarrow$ \loesung{nach rechts} \\
\es{a la izquierda} & $\rightarrow$ \loesung{nach links} \\
\es{girar} & $\rightarrow$ \loesung{abbiegen} \\
\es{al lado de} & $\rightarrow$ \loesung{neben} \\
\es{enfrente de} & $\rightarrow$ \loesung{gegenüber von} \\
\es{cerca de} & $\rightarrow$ \loesung{in der Nähe von} \\
\es{lejos de} & $\rightarrow$ \loesung{weit weg von} \\
\es{¿Dónde está...?} & $\rightarrow$ \loesung{Wo ist...?} \\
\es{¿Cómo llego a...?} & $\rightarrow$ \loesung{Wie komme ich zu...?} \\
\end{tabular}

\vspace{1em}

\textbf{Aufgabe 1: Zuordnung} (4 Punkte)

\begin{tabular}{ll}
1. todo recto & $\rightarrow$ \loesung{(3) nach links -- NEIN! (1) geradeaus} \\
\end{tabular}

Korrekte Zuordnung:
\begin{itemize}
    \item nach links = \loesung{3} (a la izquierda)
    \item in der Nähe von = \loesung{4} (cerca de)
    \item geradeaus = \loesung{1} (todo recto)
    \item nach rechts = \loesung{2} (a la derecha)
\end{itemize}

\vspace{1em}

\textbf{Aufgabe 2: Lückentext} (4 Punkte)

\begin{tabular}{ll}
a) & \loesung{enfrente de} \\
b) & \loesung{al lado de} \\
c) & \loesung{lejos} \\
d) & \loesung{todo recto} \\
\end{tabular}

\vspace{1em}

\textbf{Aufgabe 3: Dialog} (6 Punkte)

Mögliche Antworten:
\begin{itemize}
    \item Antwort 1: \loesung{Sigue todo recto. / Sigue recto por esta calle.}
    \item Antwort 2: \loesung{Gira a la derecha/izquierda. / Está al lado de...}
    \item Antwort 3: \loesung{No, está cerca. / Sí, está un poco lejos.}
\end{itemize}

\textit{Bewertung: 2P pro sinnvolle, grammatisch korrekte Antwort}

\vspace{1em}

\textbf{Aufgabe 4: Wegbeschreibung} (6 Punkte)

Beispiellösung (vom Punkt A zur Farmacia):
\begin{enumerate}
    \item \loesung{Sigue todo recto por la Calle de la Luna.}
    \item \loesung{Gira a la derecha en la Calle Mayor.}
    \item \loesung{La farmacia está a la derecha, al lado de la panadería.}
\end{enumerate}

\textit{Bewertung: 2P pro korrekten Satz mit Richtungsangabe}

% ═══════════════════════════════════════════════════════════════
% 6. HAUSAUFGABE
% ═══════════════════════════════════════════════════════════════

\section*{6. Hausaufgabe}

\begin{itemize}
    \item Lehrwerk S. 102, Aufgabe 5 (Wegbeschreibung schreiben)
    \item Vokabeln lernen: Richtungsangaben und Ortsangaben
    \item Optional: Lernpfad LP-\sequenz\block.html (digital)
\end{itemize}

% ═══════════════════════════════════════════════════════════════
% 7. REFLEXIONSNOTIZEN
% ═══════════════════════════════════════════════════════════════

\section*{7. Reflexionsnotizen (nach Durchführung)}

\begin{tabularx}{\textwidth}{|l|X|}
\hline
\textbf{Was lief gut?} & \\[2em]
\hline
\textbf{Was lief nicht?} & \\[2em]
\hline
\textbf{Zeitmanagement} & \\[2em]
\hline
\textbf{Für nächstes Mal} & \\[2em]
\hline
\end{tabularx}

\end{document}
